\documentclass[11pt,a4paper]{article}
\usepackage[utf8]{inputenc}
\usepackage[T1]{fontenc}
\usepackage[spanish,es-noshorthands]{babel}
\usepackage[margin=2.5cm]{geometry}
\usepackage{amsmath,amssymb,mathtools}
\usepackage{enumitem}
\usepackage{booktabs}
\usepackage{tikz}
\usepackage{pgfplots}
\pgfplotsset{compat=1.18}

\newcounter{ejnum}
\newenvironment{ejercicio}{\refstepcounter{ejnum}\par\medskip\noindent\textbf{Ejercicio \theejnum.}~}{\par}
\newenvironment{solucion}{\par\noindent\textit{Solución.}~}{\par\vspace{1em}}

\title{\vspace{-1.5cm}Práctica 1: Espacio muestral, sucesos y axiomas de probabilidad \\[0.3em] \large Unidad 1: Espacios de probabilidad}
\author{Probabilidad y Estadística (5765) \\ Universidad Nacional del Sur}
\date{}

\begin{document}
\maketitle

\begin{ejercicio}
Se lanza una moneda tres veces consecutivas y se registra la secuencia de caras (C) y sellos (S) obtenida.
\begin{enumerate}[label=(\alph*)]
\item Describir el espacio muestral $\Omega$.
\item Escribir por extensión el suceso $A$: ``se obtienen exactamente dos caras''.
\item Escribir por extensión el suceso $B$: ``la primera tirada es cara''.
\item Calcular $A\cap B$, $A\cup B$ y $B^c$.
\end{enumerate}
\end{ejercicio}
\begin{solucion}
\begin{enumerate}[label=(\alph*)]
\item El espacio muestral tiene $2^3=8$ resultados posibles:
\[
\Omega = \{CCC,\,CCS,\,CSC,\,SCC,\,CSS,\,SCS,\,SSC,\,SSS\}.
\]
\item $A=\{CCS,\,CSC,\,SCC\}$ (exactamente dos caras entre las tres tiradas).
\item $B=\{CCC,\,CCS,\,CSC,\,CSS\}$ (todos los resultados que comienzan con C).
\item $A\cap B = \{CCS,\,CSC\}$ (dos caras y además empieza con cara).
\newline $A\cup B = \{CCC,\,CCS,\,CSC,\,SCC,\,CSS\}$.
\newline $B^c = \{SCC,\,SCS,\,SSC,\,SSS\}$ (la primera tirada es sello).
\end{enumerate}
\end{solucion}

\begin{ejercicio}
Se selecciona al azar un artículo de una línea de producción y se clasifica según dos criterios: si tiene o no un defecto \textbf{estético} (E) y si tiene o no un defecto \textbf{funcional} (F).
\begin{enumerate}[label=(\alph*)]
\item Describir el espacio muestral en términos de los sucesos $E$ y $F$ (defecto estético / defecto funcional presentes o no).
\item Expresar en notación conjuntista el suceso ``el artículo tiene exactamente un tipo de defecto''.
\item Expresar el suceso ``el artículo no tiene ningún defecto''.
\item Expresar el suceso ``el artículo tiene al menos un defecto''.
\end{enumerate}
\end{ejercicio}
\begin{solucion}
Sea $E$ el suceso ``tiene defecto estético'' y $F$ ``tiene defecto funcional'', ambos subconjuntos del espacio muestral $\Omega$ de todos los artículos.
\begin{enumerate}[label=(\alph*)]
\item $\Omega$ se particiona en cuatro categorías posibles: $E\cap F$ (ambos defectos), $E\cap F^c$ (solo estético), $E^c\cap F$ (solo funcional), $E^c\cap F^c$ (sin defectos).
\item ``Exactamente un tipo de defecto'' $=(E\cap F^c)\cup(E^c\cap F)$ (la llamada diferencia simétrica $E \triangle F$).
\item ``Ningún defecto'' $= E^c\cap F^c = (E\cup F)^c$ (por De Morgan).
\item ``Al menos un defecto'' $= E\cup F$.
\end{enumerate}
\end{solucion}

\begin{ejercicio}
Sean $A$, $B$ y $C$ tres sucesos de un mismo espacio muestral $\Omega$. Expresar mediante operaciones de conjuntos ($\cup$, $\cap$, complemento):
\begin{enumerate}[label=(\alph*)]
\item Ocurre únicamente $A$ (ni $B$ ni $C$).
\item Ocurren exactamente dos de los tres sucesos.
\item Ocurren los tres sucesos simultáneamente.
\item No ocurre ninguno de los tres.
\end{enumerate}
\end{ejercicio}
\begin{solucion}
\begin{enumerate}[label=(\alph*)]
\item $A\cap B^c\cap C^c$.
\item $(A\cap B\cap C^c)\cup(A\cap B^c\cap C)\cup(A^c\cap B\cap C)$.
\item $A\cap B\cap C$.
\item $A^c\cap B^c\cap C^c = (A\cup B\cup C)^c$ (por De Morgan generalizado).
\end{enumerate}
\end{solucion}

\begin{ejercicio}
En un lote de 100 piezas metálicas, se detectó que 30 tienen un defecto de \textbf{dimensión} (D), 25 tienen un defecto de \textbf{acabado superficial} (S), y 10 tienen ambos defectos simultáneamente. Se elige una pieza al azar.
\begin{enumerate}[label=(\alph*)]
\item Calcular la probabilidad de que la pieza tenga al menos uno de los dos defectos.
\item Calcular la probabilidad de que la pieza no tenga ningún defecto.
\item Calcular la probabilidad de que tenga defecto de dimensión pero no de acabado.
\end{enumerate}
\end{ejercicio}
\begin{solucion}
Trabajamos con la regla de Laplace: $P(D) = 30/100=0{,}30$, $P(S)=25/100=0{,}25$, $P(D\cap S)=10/100=0{,}10$.
\begin{enumerate}[label=(\alph*)]
\item Por la regla de la suma:
\[
P(D\cup S) = P(D)+P(S)-P(D\cap S) = 0{,}30+0{,}25-0{,}10=0{,}45.
\]
\item Por complemento: $P\big((D\cup S)^c\big) = 1-0{,}45 = 0{,}55$.
\item Por diferencia: $P(D\setminus S) = P(D)-P(D\cap S) = 0{,}30-0{,}10=0{,}20$.
\end{enumerate}
\end{solucion}

\begin{ejercicio}
Demostrar, a partir de los axiomas de Kolmogorov y de las propiedades ya conocidas, que para dos sucesos $A,B\in\mathcal{A}$ cualesquiera:
\[
P(A\cap B) \geq P(A)+P(B)-1.
\]
(Esta desigualdad se conoce como \emph{cota de Bonferroni}.)
\end{ejercicio}
\begin{solucion}
Partimos de la regla de la suma, ya demostrada en la teoría:
\[
P(A\cup B) = P(A)+P(B)-P(A\cap B).
\]
Como $A\cup B \subseteq \Omega$, por el corolario de monotonía sabemos que $P(A\cup B)\leq 1$. Sustituyendo:
\[
P(A)+P(B)-P(A\cap B) \leq 1.
\]
Despejando $P(A\cap B)$ (cambiando el signo de la desigualdad al pasar de miembro):
\[
P(A\cap B) \geq P(A)+P(B)-1,
\]
como queríamos demostrar. Notar que si $P(A)+P(B)>1$, la cota es no trivial (obliga a que $A$ y $B$ se intersequen).
\end{solucion}

\begin{ejercicio}
En una empresa, el $70\%$ de los empleados usa el sistema de gestión A, el $55\%$ usa el sistema B, y se sabe que el $20\%$ no usa ninguno de los dos sistemas. Se elige un empleado al azar.
\begin{enumerate}[label=(\alph*)]
\item Calcular la probabilidad de que use ambos sistemas.
\item Calcular la probabilidad de que use exactamente uno de los dos sistemas.
\end{enumerate}
\end{ejercicio}
\begin{solucion}
Sean $A$: ``usa el sistema A'', $B$: ``usa el sistema B''. Datos: $P(A)=0{,}70$, $P(B)=0{,}55$, $P\big((A\cup B)^c\big)=0{,}20$, luego $P(A\cup B) = 1-0{,}20=0{,}80$.
\begin{enumerate}[label=(\alph*)]
\item De $P(A\cup B)=P(A)+P(B)-P(A\cap B)$:
\[
P(A\cap B) = P(A)+P(B)-P(A\cup B) = 0{,}70+0{,}55-0{,}80 = 0{,}45.
\]
\item ``Exactamente uno'' $=(A\cup B)\setminus(A\cap B)$, y como $A\cap B\subseteq A\cup B$, por monotonía:
\[
P\big((A\cup B)\setminus (A\cap B)\big) = P(A\cup B)-P(A\cap B) = 0{,}80-0{,}45=0{,}35.
\]
\end{enumerate}
\end{solucion}

\begin{ejercicio}
Un espacio muestral $\Omega$ se particiona en cuatro sucesos mutuamente excluyentes $B_1, B_2, B_3, B_4$ (un sistema completo), con probabilidades $P(B_1)=0{,}10$, $P(B_2)=0{,}20$, $P(B_3)=x$ y $P(B_4)=2x$.
\begin{enumerate}[label=(\alph*)]
\item Hallar el valor de $x$.
\item Calcular $P(B_1\cup B_3)$ y $P\big((B_2\cup B_4)^c\big)$.
\end{enumerate}
\end{ejercicio}
\begin{solucion}
\begin{enumerate}[label=(\alph*)]
\item Como $\{B_1,B_2,B_3,B_4\}$ es un sistema completo, $\sum_{i=1}^4 P(B_i)=1$:
\[
0{,}10+0{,}20+x+2x = 1 \;\Longrightarrow\; 3x = 0{,}70 \;\Longrightarrow\; x=0{,}2\overline{3}.
\]
Entonces $P(B_3)=0{,}2\overline{3}\approx 0{,}2333$ y $P(B_4)=0{,}4\overline{6}\approx0{,}4667$.
\item Como los $B_i$ son mutuamente excluyentes entre sí, por aditividad finita:
\[
P(B_1\cup B_3) = P(B_1)+P(B_3) = 0{,}10+0{,}2333=0{,}3333.
\]
\[
P(B_2\cup B_4)=P(B_2)+P(B_4)=0{,}20+0{,}4667=0{,}6667 \;\Longrightarrow\; P\big((B_2\cup B_4)^c\big)=1-0{,}6667=0{,}3333.
\]
(Coincide con el resultado anterior, ya que $(B_2\cup B_4)^c = B_1\cup B_3$ por ser $\{B_i\}$ una partición de $\Omega$.)
\end{enumerate}
\end{solucion}

\begin{ejercicio}
Se lanza un dado equilibrado de 6 caras. Consideremos los sucesos $A=\{1,2,3\}$, $B=\{2,4,6\}$ y $C=\{1,3,5\}$.
\begin{enumerate}[label=(\alph*)]
\item Verificar si $A$ y $B$ son mutuamente excluyentes.
\item Verificar si $B$ y $C$ son mutuamente excluyentes.
\item ¿Forman $B$ y $C$ un sistema completo de sucesos para $\Omega=\{1,\dots,6\}$? Justificar.
\end{enumerate}
\end{ejercicio}
\begin{solucion}
\begin{enumerate}[label=(\alph*)]
\item $A\cap B = \{1,2,3\}\cap\{2,4,6\} = \{2\} \neq \varnothing$. No son mutuamente excluyentes.
\item $B\cap C = \{2,4,6\}\cap\{1,3,5\}=\varnothing$. Sí son mutuamente excluyentes.
\item Además de ser disjuntos, $B\cup C = \{1,2,3,4,5,6\} = \Omega$, y ambos tienen probabilidad positiva ($P(B)=P(C)=1/2$). Por lo tanto $\{B,C\}$ sí es un sistema completo de sucesos: son exhaustivos y mutuamente excluyentes.
\end{enumerate}
\end{solucion}

\end{document}
